\subsubsection{Uitleg voor gebruikers}
\label{subsubsec:uitleg-voor-gebruikers}

Gebruikers worden aangespoord om waakzaam te zijn voor online fraude. Browsers en mailclients waarschuwen hen daar al langer voor. De effectiviteit van die waarschuwingen blijkt echter beperkt. \textcite{Egelman2008} toonden aan dat een groot deel van de gebruikers waarschuwingen voor phishingwebsites negeerde. Gewenning speelt hierbij een belangrijke rol, aangezien gebruikers steeds minder aandacht besteden aan waarschuwingen die ze herhaaldelijk te zien krijgen \autocite{Anderson2015}.

Naast gewenning is ook de plaatsing van een waarschuwing bepalend voor haar effectiviteit. \textcite{Petelka2019} onderzochten waarschuwingen in e-mailclients en vonden dat een waarschuwing die dicht bij de verdachte link staat, beter werkt dan een algemene melding elders in het bericht.

Naast de plaatsing speelt ook de inhoud van een waarschuwing een rol. \textcite{Egelman2008} stelden namelijk vast dat veel gebruikers de waarschuwingen niet begrepen. Technische signalen, zoals de resultaten van \emph{SPF}- of \emph{DMARC}-controles, zijn voor wie geen technische achtergrond heeft, moeilijk te duiden. Het omzetten van die signalen naar begrijpelijke taal is echter niet vanzelfsprekend, aangezien de uitleg de betekenis moet verduidelijken zonder de bewijskracht ervan te overdrijven of af te zwakken. Een uitleg moet dus in de eerste plaats begrijpelijk zijn, al garandeert dit nog geen betere beslissing. 

Onderzoek naar AI-gedreven beslissingsondersteuning toont namelijk aan dat uitleg gebruikers ertoe kan brengen het AI-systeem te volgen, ook wanneer het fout is \autocite{Bucinca2021}. Dit fenomeen wordt \emph{overreliance} genoemd. Bij phishingdetectie kan een overtuigende uitleg ervoor zorgen dat de gebruiker een foutief oordeel zal vertrouwen. 